
\documentclass[11pt]{article}
\usepackage[margin=2.5cm]{geometry}
\usepackage{graphicx}
\usepackage{booktabs}
\usepackage{float}
\usepackage{hyperref}
\usepackage{amsmath}
\usepackage{microtype}
\usepackage{url}

\title{GTEx Tissue-Expression Landscape of the RNAInterSEQ ncRNA--ncRNA Interactome}
\author{RNAInter\_sequenceParser side project}
\date{\today}

\begin{document}
\maketitle

\begin{abstract}
We link the human non-coding RNAs in the RNAInterSEQ ncRNA--ncRNA interaction
release (\texttt{interactions\_with\_sequences\_RNA\_RNA\_ncrna\_ncrna\_19-05-2026})
to their tissue expression profile in GTEx (Genotype-Tissue Expression),
which profiles postmortem human donor tissue exclusively across
54 anatomical sites. Of 10,068 unique human ncRNA entities,
7,772 (77.2\%) were resolved to at least one GENCODE gene known
to GTEx by chaining together every identifier system already present in the
dataset (Ensembl Gene ID, NCBI Gene ID via \texttt{gene2ensembl}, miRBase
mature accession via precursor-locus resolution, circBase accession via host
gene), with a GENCODE v39/v26 version fallback to recover genes pruned from
the current annotation. 6,325 of the matched entities (81.4\%)
are detected at non-zero median TPM in at least one tissue, and isoform-level
expression is available for 4,058 matched genes. Because a matched
gene's transcript set can include isoforms unrelated to the specific
molecule RNAInterSEQ has on record, and because RNAInterSEQ's own sequence
for an entity is rarely sourced from Ensembl, every candidate transcript's
sequence was independently fetched from Ensembl and tested against
RNAInterSEQ's known sequence(s) for that entity: 6,036/8,481
matched entity-locus pairs (71.2\%) have at least one such
sequence-confirmed transcript.
\end{abstract}

\section{Background}
GTEx (\url{https://gtexportal.org}) profiles RNA-seq from postmortem tissue
samples donated by human subjects. Its API v2 tissue vocabulary
(\texttt{tissueSiteDetailId}) enumerates 54 human anatomical sites
and carries no species axis -- there is no non-human GTEx data to query.
Every non-\emph{Homo sapiens} entity in the source interaction table is
therefore out of scope for this analysis by construction.

\section{Methods}

\subsection{Identifier resolution}
RNAInterSEQ stores each ncRNA entity's external identifiers in per-source
buckets (Ensembl, NCBI, miRBase, circBase, RNAcentral, GenBank, tRFdb,
``Other''). GenBank, tRFdb and ``Other'' are empty for every human entity in
this dataset and are not pursued further. The remaining four are resolved to
a GTEx-queryable gene symbol or Ensembl Gene ID as follows, tried in this
order for every entity (all are kept when more than one resolves, since a
single mature product can derive from more than one genomic locus):

\begin{enumerate}
  \item \textbf{Ensembl Gene ID} -- queried directly against GTEx
        \texttt{reference/gene}.
  \item \textbf{NCBI Gene ID} -- mapped to an Ensembl Gene ID via NCBI's own
        \texttt{gene2ensembl} table (human rows only).
  \item \textbf{Gene symbol} -- the entity's own RNAInter symbol, queried
        directly.
  \item \textbf{miRBase mature accession} (miRNA category only) -- the
        precursor locus name(s) producing that mature product are read from
        \texttt{miRNA.dat}; each precursor name is resolved to its official
        HGNC/GENCODE symbol with a single batched \texttt{mygene.info}
        alias lookup (handles irregular naming, e.g.\ \texttt{let-7a-1}
        $\rightarrow$ \texttt{MIRLET7A1}, not a fixed string transform).
  \item \textbf{circBase accession} (circRNA category only) -- circBase's
        own FASTA headers already carry the host gene symbol the circRNA
        splices from (GENCODE has no circRNA biotype, so circRNAs can only
        be matched through their linear host gene).
\end{enumerate}

\subsection{GENCODE version fallback}
Each candidate symbol/Ensembl ID is tried first at GENCODE v39 (paired with
the current \texttt{gtex\_v10} dataset), then, only for keys not found, at
v26 (paired with \texttt{gtex\_v8}). This recovers genes GENCODE has pruned
out of the current annotation -- 753 entity-locus pairs (470 of
them entities with \emph{no} v39 record at all) were only resolvable at
v26.

\subsection{Sequence-level transcript confirmation}
Matching an entity to a GENCODE gene (above) identifies the right locus, not
the right \emph{transcript}: a multi-isoform gene can carry dozens of
GENCODE transcripts, most of them unrelated to the specific molecule
RNAInterSEQ has on record. There is no shared identifier to resolve this --
RNAInterSEQ's own sequence for an entity is typically sourced from miRBase,
NCBI RefSeq, or RNAcentral, not Ensembl -- so the join key has to be the
sequence itself. For every matched gene, the spliced cDNA sequence of each
of its GENCODE transcripts is fetched from the Ensembl REST API and tested
for containment (exact match, or substring in either direction -- handling
e.g.\ a mature miRNA sequence sitting inside its precursor transcript)
against the entity's own known sequence(s).

\section{Results}

\subsection{Database composition and scope}
\begin{sloppypar}
The source file contains 119,365 ncRNA--ncRNA interaction pairs spanning
30 species-pair combinations. The dominant species are human
(81,175 pairs, 68.0\%) and mouse (35,236 pairs, 29.5\%), with the
remaining 2,954 pairs (2.5\%) covering plant, avian, rat,
and viral-host contexts (2,509 pairs involve two different species, mostly
viral miRNA--human lncRNA combinations). Because GTEx profiles exclusively
postmortem \emph{Homo sapiens} donor tissue, only the 81,175 human--human
pairs are in scope for this analysis; the full set of 57 columns produced by
this pipeline (original 33 RNAInterSEQ columns $+$ 24 GTEx annotation columns)
is stored in \path{human_ncrna_interactions_gtex_annotated.pkl}.
\end{sloppypar}

\subsection{Coverage funnel}
\begin{itemize}
  \item 10,068 unique human ncRNA entities in the 19-05-2026 dataset
  \item 7,772 (77.2\%) resolved to $\geq 1$ GENCODE gene known to GTEx
  \item 326 of those resolved to multiple genomic loci (paralogous precursors, mostly miRNA)
  \item 6,325 (81.4\% of matched) detected at median TPM $>0$ in $\geq 1$ of 54 tissues
  \item 3,332 matched genes have exactly one annotated transcript; 4,058 have multiple (49,716 unique transcripts total)
\end{itemize}

\subsection{Match rate by RNAInter category}
\begin{table}[H]
\centering
\begin{tabular}{llll}
\toprule
Category & Entities & Matched & Match rate \\
\midrule
lncRNA & 5662 & 5532 & 0.977 \\
miRNA & 3126 & 1063 & 0.340 \\
pseudo & 649 & 578 & 0.891 \\
snoRNA & 351 & 345 & 0.983 \\
circRNA & 179 & 168 & 0.939 \\
snRNA & 58 & 47 & 0.810 \\
scaRNA & 18 & 18 & 1.000 \\
rRNA & 14 & 12 & 0.857 \\
tRNA & 10 & 8 & 0.800 \\
ncRNA & 1 & 1 & 1.000 \\
\bottomrule
\end{tabular}
\caption{Entity-level GTEx match rate by RNAInter ncRNA category.}
\end{table}

\begin{figure}[H]
\centering
\includegraphics[width=0.8\textwidth]{figures/Fig1_match_rate_by_category.pdf}
\caption{GTEx match rate by RNAInter ncRNA category.}
\end{figure}

\subsection{Which identifier system resolved each match}
\begin{table}[H]
\centering
\begin{tabular}{ll}
\toprule
Resolution method & Entity-locus pairs \\
\midrule
ensembl\_id & 5275 \\
symbol & 1855 \\
mirbase\_precursor & 1086 \\
circbase\_host & 150 \\
ncbi\_to\_ensembl & 112 \\
rnacentral\_urs & 3 \\
\bottomrule
\end{tabular}
\caption{Entity-locus pairs by resolution method.}
\end{table}

\begin{figure}[H]
\centering
\includegraphics[width=0.8\textwidth]{figures/Fig4_match_method_breakdown.pdf}
\caption{Contribution of each identifier system to the final match set.}
\end{figure}

\begin{figure}[H]
\centering
\includegraphics[width=0.65\textwidth]{figures/Fig5_gencode_version_usage.pdf}
\caption{GENCODE v39/gtex\_v10 vs.\ v26/gtex\_v8 fallback usage.}
\end{figure}

RNAInter category and GTEx/GENCODE biotype agree for 79.3\% of
matched entities (primary locus per entity); the main source of disagreement
is RNAInter ``lncRNA'' entities that GENCODE annotates as protein-coding,
worth a closer look if category purity matters downstream.

\subsection{Sequence-confirmed transcripts}
6,036/8,481 matched entity-locus pairs (71.2\%)
have $\geq 1$ GENCODE transcript whose Ensembl-fetched sequence matches a
sequence RNAInterSEQ already has on record. For genes that get $\geq 1$
confirmed transcript, the confirmed set is on average 68.7\% of that
gene's total annotated transcripts (median 100.0\%) -- i.e.\ this
step meaningfully narrows multi-isoform genes down to the specific
isoform(s) RNAInterSEQ has, rather than reporting the whole gene's isoform
mixture as the entity's expression. Filtering down to only sequence-confirmed
transcripts keeps 516,186/2,698,542 (transcript,
tissue) expression rows.

\begin{table}[H]
\centering
\small
\begin{tabular}{llll}
\toprule
Category & Entity-locus pairs & Confirmed & Confirmation rate \\
\midrule
lncRNA & 5637 & 4297 & 0.762 \\
miRNA & 927 & 924 & 0.997 \\
pseudo & 591 & 391 & 0.662 \\
snoRNA & 555 & 281 & 0.506 \\
circRNA & 167 & 23 & 0.138 \\
snRNA & 164 & 82 & 0.500 \\
rRNA & 33 & 14 & 0.424 \\
scaRNA & 22 & 15 & 0.682 \\
tRNA & 8 & 8 & 1.000 \\
ncRNA & 1 & 1 & 1.000 \\
\bottomrule
\end{tabular}
\caption{Sequence-confirmation rate by RNAInter ncRNA category (entity-locus level).}
\end{table}

\begin{figure}[H]
\centering
\includegraphics[width=0.8\textwidth]{figures/Fig6_sequence_confirmation_rate.pdf}
\caption{Sequence-confirmation rate by RNAInter ncRNA category.}
\end{figure}

Confirmation rate tracks how well a category's typical sequence length and
provenance line up with GENCODE's annotated transcript boundaries: miRNA is
near-total (99.7\%, mature sequences are short and almost always found
verbatim inside their precursor transcript), while circRNA is the lowest
(13.8\%) because a circRNA's backsplice-junction sequence is, by definition,
not contiguous in the linear host transcript -- consistent with the
host-gene-proxy limitation noted below.

\subsection{Tissue specificity}
Yanai's tau across 54 GTEx tissues, computed for the 6,294 matched
genes with non-zero expression in at least one tissue: mean 0.826,
median 0.873 -- most matched ncRNAs are far from uniformly expressed.

\begin{figure}[H]
\centering
\includegraphics[width=0.7\textwidth]{figures/Fig2_tau_distribution.pdf}
\caption{Distribution of tissue specificity (tau) across matched ncRNA genes.}
\end{figure}

\begin{table}[H]
\centering
\small
\begin{tabular}{llll}
\toprule
Gene symbol & Peak tissue & Median TPM & tau \\
\midrule
MIR5004 & Testis & 0.444 & 1.000 \\
5S\_rRNA & Testis & 0.770 & 1.000 \\
LINC02203 & Testis & 1.844 & 1.000 \\
U6 & Adrenal Gland & 0.326 & 1.000 \\
C1orf167-AS1 & Testis & 3.913 & 1.000 \\
FAM27E5 & Testis & 0.508 & 1.000 \\
MGC4859 & Testis & 0.263 & 1.000 \\
LINC02632 & Testis & 0.278 & 1.000 \\
LINC01309 & Testis & 4.015 & 1.000 \\
LINC01049 & Testis & 0.309 & 1.000 \\
LINC01039 & Testis & 0.435 & 1.000 \\
LINC02530 & Testis & 0.283 & 1.000 \\
LINC01425 & Testis & 0.274 & 1.000 \\
LINC01761 & Testis & 0.050 & 1.000 \\
IQCF5-AS1 & Testis & 4.104 & 1.000 \\
\bottomrule
\end{tabular}
\caption{Top 15 most tissue-specific matched ncRNAs and their peak tissue.}
\end{table}

\subsection{Transcript multiplicity}
\begin{figure}[H]
\centering
\includegraphics[width=0.7\textwidth]{figures/Fig3_transcripts_per_gene.pdf}
\caption{Number of annotated GENCODE transcripts per matched ncRNA gene.}
\end{figure}

\section{Interaction-level tissue expression annotation}

The entity-level GTEx data are joined back to the
81,175 human--human interaction pairs in the source table
(both partners \emph{Homo sapiens}) to produce an annotated interaction
table (\path{human_ncrna_interactions_gtex_annotated.pkl}). For each
interaction row the table carries, for both partners: the set of
sequence-confirmed GENCODE transcript IDs (as a per-gene dict), tissue
specificity (tau), peak tissue, peak median TPM, and the full 54-tissue
max-TPM vector.

\begin{itemize}
  \item 54,746 (67.4\%) of pairs have both partners
        GTEx-matched.
  \item 31,264 (38.5\%) have both partners
        sequence-confirmed.
  \item 2,375 (2.9\%) share $\geq 1$ expressed GTEx tissue
        (median TPM $> 0$ in both partners).
  \item Median shared tissues across sequence-confirmed pairs:
        0 (out of 54) -- reflecting the high tissue
        specificity of ncRNAs in this dataset (median $\tau = $ 0.873)
        and the residual miRNA coverage gap.
\end{itemize}

\begin{figure}[H]
\centering
\includegraphics[width=\textwidth]{figures/Fig6_step5_funnel.pdf}
\caption{Genes and transcripts before and after sequence-confirmed
filtering (step 5), broken down by RNAInter ncRNA category.}
\end{figure}

\begin{figure}[H]
\centering
\includegraphics[width=\textwidth]{figures/Fig7_tissue_coexpression.pdf}
\caption{Number of co-expressed interaction pairs per GTEx tissue
(both partners detected at median TPM $> 0$, sequence-confirmed pairs only).}
\end{figure}

\begin{figure}[H]
\centering
\includegraphics[width=0.75\textwidth]{figures/Fig8_shared_tissue_histogram.pdf}
\caption{Distribution of the number of shared expressed tissues across
sequence-confirmed interaction pairs.}
\end{figure}

\begin{figure}[H]
\centering
\includegraphics[width=\textwidth]{figures/Fig9_category_pair_heatmap.pdf}
\caption{ncRNA category-pair co-expression landscape. Left: interaction
count (both GTEx-matched). Right: fraction of pairs with $\geq 1$ shared
expressed tissue (sequence-confirmed pairs only).}
\end{figure}

\begin{figure}[H]
\centering
\includegraphics[width=0.85\textwidth]{figures/Fig10_tau_by_category.pdf}
\caption{Yanai tissue-specificity ($\tau$) distribution per ncRNA category,
pooled across all interaction partners with GTEx data
(unique entities; $\tau = 0$ ubiquitous, $\tau = 1$ single-tissue).}
\end{figure}

\begin{figure}[H]
\centering
\includegraphics[width=0.75\textwidth]{figures/Fig11_peak_tissue_distribution.pdf}
\caption{Top 20 peak expression tissues across unique interacting ncRNA
entities (both partners pooled, deduplicated by entity).}
\end{figure}

\begin{figure}[H]
\centering
\includegraphics[width=0.65\textwidth]{figures/Fig12_tau_partner_scatter.pdf}
\caption{Tissue-specificity correlation between the two partners of each
interaction (sequence-confirmed pairs; top 5 category combinations
highlighted).}
\end{figure}

\subsection{Biological system grouping}
\begin{sloppypar}
To improve interpretability, the 54 GTEx tissue sites are collapsed into
eight biological systems: CNS (13 tissues), Peripheral nervous (2),
Cardiovascular (5), GI (10), Reproductive (9), Immune (3), Muscle (1),
and Other (11 metabolic, urological, and integumentary tissues). For each
entity, \texttt{ent1\_tissue\_systems} and \texttt{ent2\_tissue\_systems}
list the systems in which it is detected (median TPM $>0$ across confirmed
transcripts). For sequence-confirmed interaction pairs,
\texttt{shared\_tissue\_systems} gives the subset of systems shared by both
partners. Overall, 2,375 pairs (2.9\%) share expression in at
least one biological system; the most represented system is
\textbf{Reproductive} (2,168 pairs), likely driven by the broad
transcriptional activity of gonadal tissues (testis in particular ranks
among the highest expressing tissues in GTEx across all gene classes).
\end{sloppypar}

\begin{figure}[H]
\centering
\includegraphics[width=\textwidth]{figures/Fig13_tissue_system_coexpression.pdf}
\caption{Left: number of sequence-confirmed co-expressed interaction pairs
per biological system (54 GTEx tissues collapsed to 8 systems).
Right: symmetric co-occurrence matrix -- each cell shows how many pairs
have both partners expressed in system~$A$ \emph{and} system~$B$
simultaneously (a pair can contribute to multiple cells if both partners
are expressed across several systems).}
\end{figure}

\subsection{Interpreting the co-expression rate}
\begin{sloppypar}
Two questions arise naturally from the 2.9\% co-expression figure:
(\emph{i}) does the analysis account for transcripts being expressed across
multiple tissues simultaneously, and (\emph{ii}) can a transcript be
co-expressed in a tissue that GTEx simply does not detect?

\textbf{Multi-tissue expression is fully captured.} The analysis does
\emph{not} compare only peak tissues -- it uses the complete 54-tissue
max-TPM vector for each entity (columns \texttt{ent1\_tissue\_tpm} /
\texttt{ent2\_tissue\_tpm}).
\texttt{shared\_expressed\_tissues} is the intersection of the sets of tissues
where each partner has median TPM $>0$, across all confirmed transcripts and
all matched loci. Of the 31,264 sequence-confirmed pairs, 2,375
(7.6\%) share at least one tissue. Among entities with any GTEx
expression at all, the median is 39 out of 54 tissues, and
1,230 entities are expressed in every single GTEx tissue -- so when one
partner is broadly expressed the method finds overlap even if the other partner
only peaks in one tissue.

\textbf{GTEx detectability, not tissue incompatibility, drives the low rate.}
Dissecting the 31,264 sequence-confirmed pairs:
\begin{itemize}
  \item 2,375 (7.6\%) -- both partners detected, tissues overlap.
  \item 28,728 (91.9\%) -- at least one partner shows
        median TPM $= 0$ across all confirmed transcripts in all 54 GTEx tissues.
        In most cases ($\sim$68\%) the gene locus is also entirely absent from GTEx;
        in the remainder ($\sim$32\%) a different isoform of the same gene is
        detected but the sequence-confirmed isoform matching RNAInterSEQ's record is not.
        This accounts for \textbf{nearly all} non-co-expressed confirmed pairs.
  \item 161 (0.5\%) -- both partners detected in GTEx but in
        non-overlapping tissue sets within GTEx's scope.
\end{itemize}
The detectability gap is category-specific: 99.9\% of miRNA entities, 90.3\%
of snoRNA, 87.1\% of circRNA, and 100\% of tRNA show zero expression across
all 54 GTEx tissues -- for reasons inherent to how GTEx profiles RNA (see
Limitations). Only lncRNAs have reasonable coverage (31.7\% undetected).
The 2.9\% co-expression rate therefore reflects a
\textbf{conservative lower bound}: interactions involving at least one
miRNA, snoRNA, circRNA, or tRNA partner are almost never scored as
co-expressed here regardless of the underlying biology.
\end{sloppypar}

\section{Limitations}
\begin{sloppypar}
\begin{itemize}
  \item \textbf{miRNA mature/precursor mismatch.} RNAInter stores
        \emph{mature} miRNA names (e.g.\ \texttt{hsa-miR-101-3p}); GTEx
        captures bulk poly(A)-selected RNA-seq, which mostly fails to detect
        the short, rapidly processed mature miRNA at all -- even a correctly
        matched precursor (pri-miRNA) locus typically shows near-zero TPM.
        Some mature products also lack any HGNC-approved host gene symbol
        and cannot be resolved by this pipeline at all.
  \item \textbf{GTEx bulk RNA-seq averaging and isoform specificity.} GTEx
        profiles bulk RNA from heterogeneous tissue biopsies. A transcript
        expressed in a minority cell population (e.g.\ Purkinje cells within
        cerebellum, or interstitial cells within testis) may be diluted below
        detection threshold in the bulk signal and appear at TPM\,$= 0$ even
        in the tissue where it is most relevant. Additionally, the
        sequence-confirmed isoform can be non-dominant: $\sim$32\% of
        zero-tissue confirmed entities have gene-level GTEx expression in at
        least one tissue but zero TPM for the specific isoform matching
        RNAInterSEQ's sequence record. Both effects contribute to the
        dominant reason the co-expression rate is low: 28,728
        (91.9\%) of sequence-confirmed pairs have at least one
        partner with zero TPM across all confirmed transcripts in all 54 GTEx
        tissues, compared with only 161 (0.5\%) pairs where
        both partners are detected but in non-overlapping tissue sets.
  \item \textbf{GTEx tissue panel is not exhaustive.} The 54-tissue
        vocabulary covers major adult solid organs and blood but omits
        important tissue contexts where ncRNA activity has been reported:
        placenta, bone marrow, lymph node, tonsil, retina, cornea, and
        synovial tissue, among others. A co-expressed pair in any of these
        contexts will not be captured, and the \texttt{shared\_tissue\_systems}
        field will list those systems as absent for both partners.
        GTEx also captures only adult postmortem donors, missing
        developmental and foetal expression programmes.
  \item \textbf{circRNA host-gene proxy.} A circRNA's expression here is its
        \emph{linear host gene's} expression, not the circularized isoform's
        -- GENCODE has no circRNA biotype, so this is the closest available
        proxy, not a direct measurement.
  \item \textbf{RNAcentral, GenBank, tRFdb, ``Other'' buckets not pursued.}
        The latter three are empty for every human entity in this dataset.
        RNAcentral URS IDs are populated for most entities but resolving them
        to GENCODE genes would require downloading RNAcentral's full
        \texttt{id\_mapping.tsv.gz} (hundreds of MB) for marginal gains once
        Ensembl/NCBI/miRBase/circBase resolution is exhausted; left as future
        work if residual unmatched entities are specifically needed.
  \item \textbf{Multi-locus aggregation.} Entities resolved to more than one
        genomic locus (paralogous miRNA precursors, or symbol/Ensembl keys
        landing on distinct historical GENCODE records) are kept as separate
        rows rather than collapsed into one number; downstream consumers
        should decide whether to sum, average, or enumerate per locus.
  \item \textbf{Sequence confirmation is exact/containment-based, not
        alignment-based.} A transcript that differs from RNAInterSEQ's
        sequence by even one base (a different assembly patch, a SNP-bearing
        haplotype, a slightly trimmed UTR) will not be confirmed, even though
        it is almost certainly the same biological isoform -- the
        sub-100\% confirmation rates for snoRNA/snRNA/rRNA likely reflect
        this rather than a true absence of the transcript in GENCODE. A
        fuzzy/alignment-based fallback (as already used elsewhere in this
        repository, see \texttt{similarity\_computation.py}) would recover
        some of these at the cost of one pairwise alignment per candidate
        transcript.
\end{itemize}
\end{sloppypar}

\section{Data availability}
\begin{sloppypar}
All intermediate and final tables are saved as pandas pickles under
\texttt{gtex\_ncrna\_expression/data/}:
\begin{itemize}
  \item \path{ncrna_entities_human.pkl} -- the entity table built in step 1
  \item \path{gtex_gene_matches.pkl} -- the long entity-locus match table from step 2
  \item \path{gtex_gene_expression.pkl} / \path{gtex_transcript_expression.pkl} -- raw GTEx median TPM, gene and transcript level
  \item \path{ncrna_entity_locus_summary.pkl} -- the merged, analysis-ready table from step 3, including a \texttt{sequence\_confirmed} flag from step 5
  \item \path{ncrna_locus_transcript_matches.pkl} -- per (entity, gene, transcript) sequence-confirmation detail from step 5
  \item \path{gtex_transcript_expression_sequence_confirmed.pkl} -- the transcript-expression table filtered to only sequence-confirmed transcripts
  \item \path{human_ncrna_interactions_gtex_annotated.pkl} -- the 81,175 human--human interaction rows annotated with per-entity GTEx tissue expression metadata (step 6); includes \texttt{ent1\_*/ent2\_*} columns for confirmed transcripts, tau, peak tissue, the 54-tissue TPM vector, and per-entity biological system lists (\texttt{tissue\_systems}), plus \texttt{shared\_expressed\_tissues} and \texttt{shared\_tissue\_systems}
\end{itemize}
Raw GTEx/mygene.info API responses are cached under \texttt{gtex\_ncrna\_expression/cache/}.
\end{sloppypar}

\end{document}
